\documentclass[a4paper,11pt]{article}
\usepackage[utf8]{inputenc}
\usepackage[T1]{fontenc}
\usepackage{lmodern}
\usepackage[english]{babel}
\usepackage{geometry}
\geometry{left=1in, right=1in, top=1in, bottom=1in}
\usepackage{setspace}
\onehalfspacing
\usepackage{parskip}
\usepackage{enumitem}
\usepackage{booktabs}
\usepackage{graphicx}
\usepackage{amsmath}
\usepackage{natbib}
\bibliographystyle{apalike}
\usepackage[colorlinks=true, linkcolor=black, urlcolor=blue, citecolor=black]{hyperref}

\usepackage{titlesec}
\titleformat{\section}{\large\bfseries}{\thesection.}{0.5em}{}
\titleformat{\subsection}{\normalsize\bfseries}{\thesubsection}{0.5em}{}

\begin{document}

%================ TITLE PAGE ================%
\begin{center}
\vspace*{1cm}
{\LARGE \textbf{Equal Opportunity by Design?}} \\[8pt]
{\Large \textbf{The Causal Effects of Higher Education Access Policies}} \\[4pt]
{\Large \textbf{on Labour Market Inequality}} \\[8pt]
{\large Evidence from Brazil} \\[24pt]
{\large PhD Research Proposal} \\[4pt]
{\large Submitted to the Bonn Graduate School of Economics} \\[16pt]
{\large Baha Khammari} \\[4pt]
M.Sc.\ Economics, University of Cologne \\
\href{mailto:b.e.khammari@gmail.com}{b.e.khammari@gmail.com} \\[4pt]
\today
\end{center}

\vspace{24pt}

%================ INTRODUCTION ================%
\section{Introduction}

Governments worldwide rely on three principal mechanisms to expand access to higher education: scholarships, subsidised loans, and affirmative-action quotas. Each operates through a different channel---price subsidies reduce the private cost of attendance, credit instruments shift it intertemporally, and quantity mandates reserve places for underrepresented groups---and each carries distinct implications for who benefits and how. Whether these mechanisms reduce labour market inequality or merely alter the composition of graduates without changing downstream outcomes is a first-order question for education and labour policy. Yet no study to date compares their causal labour market effects within a unified empirical framework.

This gap is consequential precisely because of what recent research on perceived returns has revealed. \citet{bonevarauh2018} show that parental beliefs about the returns to educational investments differ systematically by socioeconomic status, and \citet{bonevagolin2022} find that perceived pecuniary and non-pecuniary returns explain a substantial share of the enrolment gap in postgraduate education. If the policies that expand access alter the actual returns to education---and if these actual returns differ by mechanism, intensity, and demographic group, as my master's thesis suggests---then the alignment between perceived and actual returns becomes a central question for policy design. Expanding access through scholarships may produce different actual returns than expanding it through quotas; if perceived returns do not reflect this, students may sort into the wrong programmes, and policymakers may allocate resources to the wrong instruments.

This proposal uses Brazil as an empirical laboratory for the comparison. Over a fifteen-year period, Brazil implemented all three mechanisms at national scale: \textit{ProUni} scholarships (2005), \textit{FIES} student loans (reformed 2010), and the \textit{Quota Law} (2012). The staggered rollout of each policy, combined with rich administrative employer--employee data (RAIS, approximately five million observations), creates the variation and data infrastructure needed for credible causal identification. While the empirical work is grounded in the Brazilian setting, the research questions, methods, and comparative framework are designed to inform the broader debate on how higher education access policies shape labour market inequality---a debate relevant in Germany, the EU, and other OECD countries facing similar design choices.

The project is organised as a three-paper PhD. Building on my master's thesis, which estimated the wage effects of ProUni using heterogeneity-robust staggered difference-in-differences \citep{callaway2021}, it extends the analysis in three directions: (i)~comparing scholarships against quotas and loans to identify which mechanism is most effective at reducing inequality, (ii)~examining outcomes beyond wages, including formality transitions, occupational upgrading, and job stability, and (iii)~developing a dose--response framework to estimate optimal policy intensity. Methodologically, the project contributes to the rapidly evolving literature on heterogeneous treatment effects in staggered adoption settings \citep{roth2023, dechaisemartin2020, goodmanbacon2021}.

\subsection{Motivation and relevance}

The question of whether higher education expansion translates into better labour market outcomes---and for whom---is among the most policy-relevant in both development and labour economics. Most countries use at least one of the three mechanisms studied here, yet policy choices between them are typically made without comparative causal evidence. Germany, for example, combines grant-loan hybrids (BAf\"oG), merit scholarships (Deutschlandstipendium, Begabtenf\"orderungswerke), and an ongoing debate about social quotas in university admissions. The EU more broadly grapples with cross-national differences in tuition, financial aid, and equity mandates under the Bologna Process. Evidence on which mechanism works best---and for whom---is directly transferable across these settings.

Brazil provides a uniquely clean setting for this comparison. Tertiary enrolment more than doubled between 2003 and 2018, driven primarily by the three policies studied here. This expansion occurred alongside a historic reduction in income inequality (the Gini coefficient fell from 0.59 to 0.52 between 2001 and 2014), but the contribution of education policy to this trend is causally unidentified.

My thesis findings motivate the core questions. ProUni's wage effects are non-monotonic (+3.04\% in low-dose regions, near-zero in high-dose regions), gains accrue disproportionately to non-white workers, and na\"ive two-way fixed effects (TWFE) estimation masks these patterns entirely. These results raise three questions that this PhD project addresses: Do loans and quotas produce similarly heterogeneous effects? Are scholarship-based and quota-based mechanisms substitutes or complements? And is there an optimal intensity that maximises returns for disadvantaged groups?

%================ BACKGROUND ================%
\section{Background and Scientific Basis}

\subsection{Institutional context: Brazil as empirical setting}

The comparative analysis requires a setting where multiple access mechanisms were introduced at scale with variation in timing and intensity. Brazil meets these conditions.

\textbf{ProUni} (\textit{Programa Universidade para Todos}, Lei n\textsuperscript{o} 11.096/2005) provides full or partial tuition scholarships at private higher education institutions to students from low-income households (family income $\leq 1.5$ or $\leq 3$ minimum wages per capita, respectively) who attended public secondary schools or received full private-school scholarships. Institutions participate voluntarily in exchange for tax exemptions. By 2019, ProUni had awarded over 2.4 million scholarships.

\textbf{FIES} (\textit{Fundo de Financiamento Estudantil}) is a federal student loan programme for private-institution students. Following a major expansion in 2010, FIES disbursed over 2.1 million loans by 2014, with subsidised interest rates and income-contingent repayment. A fiscal consolidation in 2015 sharply reduced new contracts, creating temporal variation useful for causal identification.

\textbf{The Quota Law} (Lei n\textsuperscript{o} 12.711/2012) mandates that federal universities reserve at least 50\% of places for graduates of public secondary schools, with sub-quotas for low-income students and ethnic minorities (Black, Brown/\textit{pardo}, and Indigenous). Implementation was staggered: universities had until 2016 to reach full compliance, creating the variation in treatment timing required for difference-in-differences estimation.

\subsection{Gaps in the literature}

The broader literature on higher education access and labour markets evaluates individual policy instruments in isolation---scholarships in one study, quotas in another, loans in a third---without comparing mechanisms head-to-head. This fragmentation limits what policymakers can learn: knowing that a scholarship programme raises wages does not tell them whether the same resources allocated to quotas or loans would have been more effective. The gap is not specific to Brazil; it characterises the literature globally.

A parallel literature, to which Prof.\ Boneva has made central contributions, studies the \textit{demand} side: why students and families choose to invest in education, and how beliefs about returns shape these decisions. \citet{bonevarauh2018} show that parents perceive higher returns to late (relative to early) educational investments, and that perceived returns to early investment correlate positively with household income---a pattern that can generate intergenerational persistence in educational inequality. \citet{bonevaetal2022} document that perceived returns to different types of investments (time, materials, school quality) predict actual investment behaviour. \citet{bonevagolin2022} find that perceived pecuniary and non-pecuniary returns explain a substantial share of the socioeconomic gap in postgraduate enrolment intentions, with the perceived consumption value of education accounting for roughly half the variation.

These findings establish that beliefs about returns matter for educational decisions. What is missing is the supply-side counterpart: rigorous estimates of \textit{actual} causal returns, differentiated by access mechanism, for the populations these policies target. If actual returns are heterogeneous---varying by policy instrument, dosage, gender, and ethnicity, as my thesis results suggest---then the gap between perceived and actual returns may differ systematically across mechanisms and demographic groups. This project provides those estimates.

Within the Brazilian case, studies of ProUni focus on enrolment effects \citep{cataldo2020} and institutional quality \citep{aprile2009}, with limited attention to downstream labour market outcomes. Research on the Quota Law examines academic performance and graduation rates but rarely traces beneficiaries into the formal labour market. FIES has received attention primarily through its fiscal costs and default rates.

Methodologically, most existing evaluations rely on conventional TWFE estimators, which are biased under staggered adoption with heterogeneous treatment effects \citep{goodmanbacon2021, dechaisemartin2020}. My thesis demonstrated this empirically for ProUni: the TWFE estimate of $-0.65\%$ on wages reverses to $+3.04\%$ in low-dose regions once heterogeneity-robust methods are applied. Similar biases likely affect existing estimates for FIES, the Quota Law, and education access policies evaluated in other countries.

%================ RESEARCH QUESTIONS ================%
\section{Research Questions and Challenges}

\subsection{Research questions}

The project is organised around three self-contained papers, each addressing a distinct research question:

\textbf{Paper~1: Scholarships vs.\ Quotas---Comparative Labour Market Returns.}
\begin{quote}
\textit{Do ProUni scholarships and Quota Law reservations produce different effects on formal-sector wages, employment, and occupational attainment, and do these effects vary by gender, ethnicity, and region?}
\end{quote}
This paper extends my thesis framework to the Quota Law and provides the first head-to-head comparison. It exploits the staggered rollout of quota compliance across federal universities (2013--2016) and the regional variation in ProUni scholarship allocation. The comparison is of particular interest in light of \citeauthor{bonevarauh2017}'s (\citeyear{bonevarauh2017}) finding that socioeconomic enrolment gaps are driven by both pecuniary and non-pecuniary perceived returns: scholarships and quotas differ in precisely these dimensions (scholarships reduce costs; quotas guarantee access without changing prices), so comparing their labour market returns can reveal which channel matters more for downstream outcomes.

\textbf{Paper~2: Beyond Wages---Formality, Occupational Upgrading, and Job Stability.}
\begin{quote}
\textit{Do higher education access policies affect labour market outcomes beyond wages, including formal-sector entry, occupational mobility, and employment duration?}
\end{quote}
This paper broadens the outcome set. Brazil's large informal sector (approximately 40\% of employment) means that wage effects in the formal sector capture only part of the return to education. RAIS records allow tracking of formal-sector entry, job tenure, occupational codes (CBO), and employer characteristics over time. These non-wage outcomes correspond to the non-pecuniary returns that \citet{bonevarauh2017} identify as a major driver of enrolment decisions---job stability, occupational prestige, and formality are dimensions students and families likely value beyond their effect on wages.

\textbf{Paper~3: Dose--Response and Optimal Policy Design.}
\begin{quote}
\textit{What is the shape of the dose--response relationship between higher education access intensity and labour market outcomes, and what does it imply for optimal policy calibration?}
\end{quote}
This paper builds on my thesis finding of non-monotonic ProUni effects. It formalises the dose--response estimation across all three policies using continuous treatment frameworks \citep{callaway2024} and asks whether there is an identifiable ``saturation point'' beyond which additional access yields diminishing or zero returns.

\subsection{Anticipated challenges}

\begin{itemize}[nosep]
    \item \textbf{Selection into treatment.} ProUni and FIES are voluntary; the Quota Law mandates institutional behaviour but not individual take-up. Each paper must address selection at the appropriate level (microregion, institution, or individual) using doubly-robust methods and sensitivity analysis.
    \item \textbf{Concurrent policies.} The three programmes overlap temporally and geographically. Isolating each policy's contribution requires careful control for the others, using differences in eligibility criteria and rollout timing.
    \item \textbf{Data linkage.} Individual-level linkage across programme records and RAIS is not feasible without CPF identifiers. The project therefore operates at the microregion--year level, following the design of my thesis, supplemented by institution-level analysis where data permit.
    \item \textbf{RAIS coverage.} RAIS captures only formal employment. Effects on informality must be inferred indirectly, e.g.\ through PNAD Cont\'inua survey data or changes in the formal employment share at the microregion level.
\end{itemize}

%================ METHODOLOGY ================%
\section{Methodology}

\subsection{Data sources}

All data are publicly accessible through the \texttt{basedosdados} BigQuery interface, which I used extensively in my thesis to construct a microregion--year panel of 558 regions over 2005--2019.

\begin{table}[h]
\centering
\small
\begin{tabular}{@{}lll@{}}
\toprule
\textbf{Dataset} & \textbf{Source} & \textbf{Key variables} \\
\midrule
RAIS & Ministry of Labour & Wages, occupation, tenure, sector, employer \\
ProUni & Ministry of Education & Scholarships by institution, type, year \\
FIES & FNDE & Loans by institution, amount, year \\
Quota Law compliance & Federal universities & Reservation shares, implementation year \\
Census / PNAD Cont\'inua & IBGE & Demographics, education, informality \\
\bottomrule
\end{tabular}
\end{table}

\subsection{Identification strategy}

\textbf{Papers 1 and 2} use the \citet{callaway2021} group--time average treatment effect estimator, which avoids the bias from ``forbidden comparisons'' in staggered TWFE by restricting contrasts to not-yet-treated or never-treated units. Treatment is defined at the microregion level as the year in which scholarship/quota exposure first exceeds a threshold, following the approach validated in my thesis. Doubly-robust estimation combines inverse probability weighting with outcome regression to reduce sensitivity to model misspecification.

For the Quota Law specifically, the staggered compliance timeline (2013--2016) across 63 federal universities provides clean variation in treatment timing. I construct a municipality-level measure of quota exposure based on the share of the local student-age population within commuting distance of a treated university.

\textbf{Paper~3} extends the binary treatment framework to continuous treatment intensity using the generalised DiD estimator for continuous treatments \citep{callaway2024}. This allows nonparametric estimation of the dose--response function $\tau(d)$, where $d$ measures cumulative scholarship/quota intensity per capita.

\subsection{Robustness and sensitivity}

Following the protocol established in my thesis, each paper includes:
\begin{itemize}[nosep]
    \item Pre-trend testing and HonestDiD sensitivity analysis \citep{rambachan2023} under violations of parallel trends.
    \item Bacon decomposition \citep{goodmanbacon2021} to diagnose TWFE bias.
    \item Placebo tests using pre-treatment cohorts and untreated regions.
    \item Controls for concurrent policies (e.g., \textit{Bolsa Fam\'ilia}, minimum wage adjustments).
    \item Alternative treatment definitions and dose cutoffs.
\end{itemize}

\subsection{Software}

All estimation is conducted in R, using the \texttt{did} package \citep{callaway2021}, \texttt{fixest} for high-dimensional fixed effects \citep{berge2018}, \texttt{HonestDiD} \citep{rambachan2023}, and \texttt{tidyverse}/\texttt{ggplot2} for data processing and visualisation. Data extraction from BigQuery uses SQL via \texttt{basedosdados}.

%================ EXPECTED CONTRIBUTION ================%
\section{Expected Contribution and Impact}

\subsection{Scientific contribution}

\begin{enumerate}[nosep]
    \item \textbf{First head-to-head comparison} of the three main higher education access mechanisms (scholarships, loans, quotas) within a unified causal framework, addressing a gap in the education and labour economics literatures that extends well beyond the Brazilian case.
    \item \textbf{Supply-side complement} to the perceived-returns literature. While \citet{bonevarauh2018} and \citet{bonevagolin2022} establish that beliefs about returns drive enrolment decisions, this project provides the causal estimates of \textit{actual} returns by mechanism, dosage, and demographic group needed to assess whether those beliefs are well-calibrated---and where they are not.
    \item \textbf{Expanded outcome analysis} beyond wages to formality, occupational upgrading, and job stability---outcomes that are first-order relevant wherever informal employment is prevalent but rarely examined in the education policy literature.
    \item \textbf{Dose--response estimation} for education access policies, providing evidence on optimal programme intensity and informing the calibration of future interventions.
    \item \textbf{Applied implementation} of heterogeneity-robust DiD methods in a staggered multi-policy environment, building on \citet{callaway2021}, \citet{dechaisemartin2020}, and \citet{sun2021}, with a comparative structure that serves as a template for other country settings.
\end{enumerate}

\subsection{Policy impact}

The findings speak to any country choosing between scholarships, loans, and quotas as instruments for widening higher education access. In Brazil, the Quota Law's ten-year review (mandated for 2022, ongoing) and recurrent debates about FIES's fiscal sustainability create immediate demand for rigorous comparative evidence. The dose--response analysis in Paper~3 can inform decisions about programme scale---a question my thesis flagged (diminishing returns in high-dose regions) but could not fully resolve with a single-policy design.

\subsection{External validity and generalisability}

The research design is grounded in Brazil, but the underlying question---which access mechanism reduces labour market inequality most effectively---is universal. The three mechanisms studied here (scholarships, loans, quotas) map directly onto policy instruments used in other contexts:

\begin{itemize}[nosep]
    \item \textbf{Germany}: BAf\"oG (grant-loan hybrid), the Deutschlandstipendium, Begabtenf\"orderungswerke, and recent proposals for social quotas in university admissions.
    \item \textbf{EU}: Cross-national variation in tuition regimes, means-tested grants, income-contingent loans, and equity mandates under the Bologna Process.
    \item \textbf{Global South}: India's reservation system, South Africa's NSFAS loans, and scholarship programmes across Latin America and Southeast Asia.
\end{itemize}

The methodological framework---heterogeneity-robust DiD under staggered adoption---transfers directly to any setting with phased policy rollout. The comparative structure (which mechanism, for whom, at what intensity) provides a template that can be applied to other countries' higher education systems as administrative data become available.

%================ FIT ================%
\section{Fit with the BGSE and Proposed Supervision}

This project fits the research environment at Bonn on several dimensions. Prof.\ Boneva's work on perceived returns to education and their role in enrolment decisions \citep{bonevarauh2018, bonevagolin2022} addresses the demand side of the question I study from the supply side. Her finding that non-pecuniary returns---the consumption value of education, job satisfaction, social status---explain much of the socioeconomic enrolment gap \citep{bonevarauh2017} connects directly to Paper~2, which extends outcome measurement beyond wages to formality, occupational upgrading, and job stability: the non-wage dimensions of labour market returns that students and families likely value. The gender asymmetry I document in my thesis---persistent wage gaps despite majority-female scholarship uptake---connects to her ERC Starting Grant (BELIEFS) on gender inequality in labour markets and the mechanisms behind child penalties.

Methodologically, Prof.\ Boneva's approach (large-scale survey data, belief elicitation, RCTs) and mine (administrative data, difference-in-differences) are complementary. A natural extension of the PhD project would combine both: using my causal estimates of actual returns by mechanism with survey data on perceived returns to quantify the beliefs--reality gap and its consequences for enrolment decisions and policy design. This would connect the two research agendas in a way that neither administrative data nor survey data alone can achieve.

More broadly, the ECONtribute Cluster of Excellence, with its focus on markets and public policy, and the BGSE's partnership with IZA provide an environment where the labour economics and policy evaluation dimensions of this project are well supported. Prof.\ Schiprowski's research on labour market policy evaluation using administrative data and Prof.\ von Gaudecker's applied microeconometric work on wage inequality offer additional methodological and substantive overlap.

I see this project as the first step in a broader research agenda on the design of equity-oriented education policy. The methods and comparative framework developed here extend naturally to related topics---the labour market effects of vocational training, apprenticeship systems, or lifelong learning policies---in Germany and elsewhere.

%================ RESEARCH ETHICS ================%
\section{Research Ethics}

This project uses exclusively \textbf{publicly available, de-identified administrative data}. RAIS microdata accessed through \texttt{basedosdados} are aggregated at the establishment or microregion level; individual worker identifiers (CPF) are not available in the public release. ProUni and FIES records are published as aggregate statistics by institution and year. Census and PNAD data are anonymised by IBGE prior to release.

No human subjects are directly involved. The project does not require institutional review board (IRB/ethics committee) approval under standard guidelines for secondary analysis of public administrative data, though I will seek confirmation from the host institution's ethics office upon enrolment.

The research addresses inequality by design: it examines whether affirmative-action and access policies achieve their stated equity objectives. Findings will be presented with care to avoid stigmatising specific demographic groups. Following best practice in the economics of discrimination, I report heterogeneous effects by ethnicity, gender, and region to illuminate policy-relevant disparities, not to attribute outcomes to group characteristics.

All code and replication materials will be made publicly available on GitHub upon publication.

%================ TIMELINE ================%
\section{Timeline}

The proposed timeline assumes a structured PhD programme at the BGSE beginning in October 2026. The first two years of the BGSE programme combine coursework with research; the timeline accounts for this by placing the main estimation work in Year~1 (months 5--12) and Year~2, after core coursework is completed.

\begin{table}[h]
\centering
\small
\begin{tabular}{@{}lp{10.5cm}@{}}
\toprule
\textbf{Period} & \textbf{Milestones} \\
\midrule
\textbf{Year 1} & \\
Months 1--4 & BGSE coursework; literature review. Extend thesis data infrastructure to include FIES and Quota Law records. \\
Months 5--8 & \textbf{Paper 1} estimation: Quota Law treatment assignment, staggered DiD, comparative analysis with ProUni results. \\
Months 9--12 & Paper 1 draft completed. Present at internal seminar. Begin Paper 2 data construction (occupational codes, tenure, formality measures). \\[4pt]
\textbf{Year 2} & \\
Months 13--16 & \textbf{Paper 2} estimation: formality transitions, occupational upgrading, job stability. \\
Months 17--20 & Paper 2 draft completed. Submit Paper 1 to journal. Conference presentations (e.g., EEA, VfS, ESPE, LACEA). \\
Months 21--24 & Begin \textbf{Paper 3}: continuous treatment framework, dose--response estimation. Revise Paper 1 based on feedback. \\[4pt]
\textbf{Year 3} & \\
Months 25--28 & Paper 3 estimation and draft. Submit Paper 2. \\
Months 29--32 & Revisions and integration. Job market preparation. \\
Months 33--36 & Dissertation submission. Job market presentations. \\
\bottomrule
\end{tabular}
\end{table}

%================ REFERENCES ================%
\section{References}

\begingroup
\renewcommand{\section}[2]{}
\begin{thebibliography}{99}

\bibitem[Aprile and Barone, 2009]{aprile2009}
Aprile, M.R. and Barone, R.E.M. (2009).
\newblock Educa\c{c}\~ao superior: pol\'iticas p\'ublicas para inclus\~ao social.
\newblock \textit{Revista @mbienteeduca\c{c}\~ao}, 2(1), 39--55.

\bibitem[Berg\'e, 2018]{berge2018}
Berg\'e, L. (2018).
\newblock Efficient estimation of maximum likelihood models with multiple fixed-effects: the R package fixest.
\newblock \textit{CREA Discussion Paper}, 2018-13.

\bibitem[Boneva and Rauh, 2017]{bonevarauh2017}
Boneva, T. and Rauh, C. (2017).
\newblock Socio-economic gaps in university enrollment: the role of perceived pecuniary and non-pecuniary returns.
\newblock \textit{CESifo Working Paper}, No.\ 6756.

\bibitem[Boneva and Rauh, 2018]{bonevarauh2018}
Boneva, T. and Rauh, C. (2018).
\newblock Parental beliefs about returns to educational investments---the later the better?
\newblock \textit{Journal of the European Economic Association}, 16(6), 1669--1712.

\bibitem[Boneva, Attanasio, and Rauh, 2022]{bonevaetal2022}
Boneva, T., Attanasio, O. and Rauh, C. (2022).
\newblock Parental beliefs about returns to different types of investments in school children.
\newblock \textit{Journal of Human Resources}, 57(6), 1789--1825.

\bibitem[Boneva, Golin, and Rauh, 2022]{bonevagolin2022}
Boneva, T., Golin, M. and Rauh, C. (2022).
\newblock Can perceived returns explain enrollment gaps in postgraduate education?
\newblock \textit{Labour Economics}, 77, 101998.

\bibitem[Callaway and Sant'Anna, 2021]{callaway2021}
Callaway, B. and Sant'Anna, P.H.C. (2021).
\newblock Difference-in-differences with multiple time periods.
\newblock \textit{Journal of Econometrics}, 225(2), 200--230.

\bibitem[Callaway et~al., 2024]{callaway2024}
Callaway, B., Goodman-Bacon, A. and Sant'Anna, P.H.C. (2024).
\newblock Difference-in-differences with a continuous treatment.
\newblock \textit{NBER Working Paper}, No.\ 32117.

\bibitem[Cataldo et~al., 2020]{cataldo2020}
Cataldo, B., Cede\~no, D.K.A. and Di~Maio, M. (2020).
\newblock Scholarships and student outcomes: evidence from Brazil's ProUni.
\newblock \textit{IZA Discussion Paper}, No.\ 13748.

\bibitem[de~Chaisemartin and D'Haultf\oe ille, 2020]{dechaisemartin2020}
de~Chaisemartin, C. and D'Haultf\oe ille, X. (2020).
\newblock Two-way fixed effects estimators with heterogeneous treatment effects.
\newblock \textit{American Economic Review}, 110(9), 2964--2996.

\bibitem[Goodman-Bacon, 2021]{goodmanbacon2021}
Goodman-Bacon, A. (2021).
\newblock Difference-in-differences with variation in treatment timing.
\newblock \textit{Journal of Econometrics}, 225(2), 254--277.

\bibitem[Rambachan and Roth, 2023]{rambachan2023}
Rambachan, A. and Roth, J. (2023).
\newblock A more credible approach to parallel trends.
\newblock \textit{Review of Economic Studies}, 90(5), 2555--2591.

\bibitem[Roth et~al., 2023]{roth2023}
Roth, J., Sant'Anna, P.H.C., Bilinski, A. and Poe, J. (2023).
\newblock What's trending in difference-in-differences? A synthesis of the recent econometrics literature.
\newblock \textit{Journal of Econometrics}, 235(2), 2218--2244.

\bibitem[Sun and Abraham, 2021]{sun2021}
Sun, L. and Abraham, S. (2021).
\newblock Estimating dynamic treatment effects in event studies with heterogeneous treatment effects.
\newblock \textit{Journal of Econometrics}, 225(2), 175--199.

\end{thebibliography}
\endgroup

\end{document}
